\section{Gaussian Space --- Condorcet}
%%% generated by the blueprint pipeline from TCSlib.BooleanAnalysis.GaussianSpace.Condorcet
\subsection*{Gaussian Space --- Condorcet: Introduction}

Declarations of the Lean module \texttt{TCSlib.BooleanAnalysis.GaussianSpace.Condorcet}.

\subsection*{Gaussian Space --- Condorcet: Results}

\begin{theorem}[Small influences and Condorcet rationality]
\lean{BooleanAnalysis.GaussianSpace.condorcet_rationality_small_influences}
\label{BooleanAnalysis.GaussianSpace.condorcet_rationality_small_influences}
\leanok
\uses{ArrowTheorem.Profile, ArrowTheorem.abVotes, ArrowTheorem.bcVotes, ArrowTheorem.caVotes, BooleanAnalysis.BooleanFunc, BooleanAnalysis.ThresholdFunctions.maxInfluence, BooleanAnalysis.isPmOne}
For every $\eta>0$, there is a $\tau>0$ such that, for every natural number $n$ and
every function $f:\{0,1\}^n\to\mathbb R$ taking only the values $-1$ and $1$, if its
largest absolute coordinate influence is at most $\tau$, then the fraction of profiles
assigning each of $n$ voters one of the six strict rankings of three alternatives for
which the three pairwise social outputs are not all equal is at most
$\frac34+\frac{3}{2\pi}\arcsin(\frac13)+\eta$.
\end{theorem}

\begin{theorem}[Asymptotic majority rationality probability]
\lean{BooleanAnalysis.GaussianSpace.majority_condorcet_rationality_tendsto}
\label{BooleanAnalysis.GaussianSpace.majority_condorcet_rationality_tendsto}
\leanok
\uses{ArrowTheorem.Profile, ArrowTheorem.abVotes, ArrowTheorem.bcVotes, ArrowTheorem.caVotes, BooleanAnalysis.majority}
For each $k\ge 0$, choose a profile of $2k+1$ voters independently and uniformly from
the six strict rankings of three alternatives. Each of the three pairwise elections is
decided by majority. As $k\to\infty$, the probability that the three majority outcomes
are not all equal converges to $\frac34+\frac{3}{2\pi}\arcsin\!\left(\frac13\right)$.
\end{theorem}
